\IfFileExists{stacks-project.cls}{%
\documentclass{stacks-project}
}{%
\documentclass{amsart}
}

% For dealing with references we use the comment environment
\usepackage{verbatim}
\newenvironment{reference}{\comment}{\endcomment}
%\newenvironment{reference}{}{}
\newenvironment{slogan}{\comment}{\endcomment}
\newenvironment{history}{\comment}{\endcomment}

% For commutative diagrams we use Xy-pic
\usepackage[all]{xy}

% We use 2cell for 2-commutative diagrams.
\xyoption{2cell}
\UseAllTwocells

% We use multicol for the list of chapters between chapters
\usepackage{multicol}

% This is generally recommended for better output
\usepackage{lmodern}
\usepackage[T1]{fontenc}

% For cross-file-references
\usepackage{xr-hyper}

% Package for hypertext links:
\usepackage{hyperref}

% For any local file, say "hello.tex" you want to link to please
% use \externaldocument[hello-]{hello}
\externaldocument[introduction-]{introduction}
\externaldocument[conventions-]{conventions}
\externaldocument[sets-]{sets}
\externaldocument[categories-]{categories}
\externaldocument[topology-]{topology}
\externaldocument[sheaves-]{sheaves}
\externaldocument[sites-]{sites}
\externaldocument[stacks-]{stacks}
\externaldocument[fields-]{fields}
\externaldocument[algebra-]{algebra}
\externaldocument[brauer-]{brauer}
\externaldocument[homology-]{homology}
\externaldocument[derived-]{derived}
\externaldocument[simplicial-]{simplicial}
\externaldocument[more-algebra-]{more-algebra}
\externaldocument[smoothing-]{smoothing}
\externaldocument[modules-]{modules}
\externaldocument[sites-modules-]{sites-modules}
\externaldocument[injectives-]{injectives}
\externaldocument[cohomology-]{cohomology}
\externaldocument[sites-cohomology-]{sites-cohomology}
\externaldocument[dga-]{dga}
\externaldocument[dpa-]{dpa}
\externaldocument[sdga-]{sdga}
\externaldocument[hypercovering-]{hypercovering}
\externaldocument[schemes-]{schemes}
\externaldocument[constructions-]{constructions}
\externaldocument[properties-]{properties}
\externaldocument[morphisms-]{morphisms}
\externaldocument[coherent-]{coherent}
\externaldocument[divisors-]{divisors}
\externaldocument[limits-]{limits}
\externaldocument[varieties-]{varieties}
\externaldocument[topologies-]{topologies}
\externaldocument[descent-]{descent}
\externaldocument[perfect-]{perfect}
\externaldocument[more-morphisms-]{more-morphisms}
\externaldocument[flat-]{flat}
\externaldocument[groupoids-]{groupoids}
\externaldocument[more-groupoids-]{more-groupoids}
\externaldocument[etale-]{etale}
\externaldocument[chow-]{chow}
\externaldocument[intersection-]{intersection}
\externaldocument[pic-]{pic}
\externaldocument[weil-]{weil}
\externaldocument[adequate-]{adequate}
\externaldocument[dualizing-]{dualizing}
\externaldocument[duality-]{duality}
\externaldocument[discriminant-]{discriminant}
\externaldocument[derham-]{derham}
\externaldocument[local-cohomology-]{local-cohomology}
\externaldocument[algebraization-]{algebraization}
\externaldocument[curves-]{curves}
\externaldocument[resolve-]{resolve}
\externaldocument[models-]{models}
\externaldocument[functors-]{functors}
\externaldocument[equiv-]{equiv}
\externaldocument[pione-]{pione}
\externaldocument[etale-cohomology-]{etale-cohomology}
\externaldocument[proetale-]{proetale}
\externaldocument[relative-cycles-]{relative-cycles}
\externaldocument[more-etale-]{more-etale}
\externaldocument[trace-]{trace}
\externaldocument[crystalline-]{crystalline}
\externaldocument[spaces-]{spaces}
\externaldocument[spaces-properties-]{spaces-properties}
\externaldocument[spaces-morphisms-]{spaces-morphisms}
\externaldocument[decent-spaces-]{decent-spaces}
\externaldocument[spaces-cohomology-]{spaces-cohomology}
\externaldocument[spaces-limits-]{spaces-limits}
\externaldocument[spaces-divisors-]{spaces-divisors}
\externaldocument[spaces-over-fields-]{spaces-over-fields}
\externaldocument[spaces-topologies-]{spaces-topologies}
\externaldocument[spaces-descent-]{spaces-descent}
\externaldocument[spaces-perfect-]{spaces-perfect}
\externaldocument[spaces-more-morphisms-]{spaces-more-morphisms}
\externaldocument[spaces-flat-]{spaces-flat}
\externaldocument[spaces-groupoids-]{spaces-groupoids}
\externaldocument[spaces-more-groupoids-]{spaces-more-groupoids}
\externaldocument[bootstrap-]{bootstrap}
\externaldocument[spaces-pushouts-]{spaces-pushouts}
\externaldocument[spaces-chow-]{spaces-chow}
\externaldocument[groupoids-quotients-]{groupoids-quotients}
\externaldocument[spaces-more-cohomology-]{spaces-more-cohomology}
\externaldocument[spaces-simplicial-]{spaces-simplicial}
\externaldocument[spaces-duality-]{spaces-duality}
\externaldocument[formal-spaces-]{formal-spaces}
\externaldocument[restricted-]{restricted}
\externaldocument[spaces-resolve-]{spaces-resolve}
\externaldocument[formal-defos-]{formal-defos}
\externaldocument[defos-]{defos}
\externaldocument[cotangent-]{cotangent}
\externaldocument[examples-defos-]{examples-defos}
\externaldocument[algebraic-]{algebraic}
\externaldocument[examples-stacks-]{examples-stacks}
\externaldocument[stacks-sheaves-]{stacks-sheaves}
\externaldocument[criteria-]{criteria}
\externaldocument[artin-]{artin}
\externaldocument[quot-]{quot}
\externaldocument[stacks-properties-]{stacks-properties}
\externaldocument[stacks-morphisms-]{stacks-morphisms}
\externaldocument[stacks-limits-]{stacks-limits}
\externaldocument[stacks-cohomology-]{stacks-cohomology}
\externaldocument[stacks-perfect-]{stacks-perfect}
\externaldocument[stacks-introduction-]{stacks-introduction}
\externaldocument[stacks-more-morphisms-]{stacks-more-morphisms}
\externaldocument[stacks-geometry-]{stacks-geometry}
\externaldocument[moduli-]{moduli}
\externaldocument[moduli-curves-]{moduli-curves}
\externaldocument[examples-]{examples}
\externaldocument[exercises-]{exercises}
\externaldocument[guide-]{guide}
\externaldocument[desirables-]{desirables}
\externaldocument[coding-]{coding}
\externaldocument[obsolete-]{obsolete}
\externaldocument[fdl-]{fdl}
\externaldocument[index-]{index}

% Theorem environments.
%
\theoremstyle{plain}
\newtheorem{theorem}[subsection]{Theorem}
\newtheorem{proposition}[subsection]{Proposition}
\newtheorem{lemma}[subsection]{Lemma}

\theoremstyle{definition}
\newtheorem{definition}[subsection]{Definition}
\newtheorem{example}[subsection]{Example}
\newtheorem{exercise}[subsection]{Exercise}
\newtheorem{situation}[subsection]{Situation}

\theoremstyle{remark}
\newtheorem{remark}[subsection]{Remark}
\newtheorem{remarks}[subsection]{Remarks}

\numberwithin{equation}{subsection}

% Macros
%
\def\lim{\mathop{\mathrm{lim}}\nolimits}
\def\colim{\mathop{\mathrm{colim}}\nolimits}
\def\Spec{\mathop{\mathrm{Spec}}}
\def\Hom{\mathop{\mathrm{Hom}}\nolimits}
\def\Ext{\mathop{\mathrm{Ext}}\nolimits}
\def\SheafHom{\mathop{\mathcal{H}\!\mathit{om}}\nolimits}
\def\SheafExt{\mathop{\mathcal{E}\!\mathit{xt}}\nolimits}
\def\Sch{\mathit{Sch}}
\def\Mor{\mathop{\mathrm{Mor}}\nolimits}
\def\Ob{\mathop{\mathrm{Ob}}\nolimits}
\def\Sh{\mathop{\mathit{Sh}}\nolimits}
\def\NL{\mathop{N\!L}\nolimits}
\def\CH{\mathop{\mathrm{CH}}\nolimits}
\def\proetale{{pro\text{-}\acute{e}tale}}
\def\etale{{\acute{e}tale}}
\def\QCoh{\mathit{QCoh}}
\def\Ker{\mathop{\mathrm{Ker}}}
\def\Im{\mathop{\mathrm{Im}}}
\def\Coker{\mathop{\mathrm{Coker}}}
\def\Coim{\mathop{\mathrm{Coim}}}

% Boxtimes
%
\DeclareMathSymbol{\boxtimes}{\mathbin}{AMSa}{"02}

%
% Macros for moduli stacks/spaces
%
\def\QCohstack{\mathcal{QC}\!\mathit{oh}}
\def\Cohstack{\mathcal{C}\!\mathit{oh}}
\def\Spacesstack{\mathcal{S}\!\mathit{paces}}
\def\Quotfunctor{\mathrm{Quot}}
\def\Hilbfunctor{\mathrm{Hilb}}
\def\Curvesstack{\mathcal{C}\!\mathit{urves}}
\def\Polarizedstack{\mathcal{P}\!\mathit{olarized}}
\def\Complexesstack{\mathcal{C}\!\mathit{omplexes}}
% \Pic is the operator that assigns to X its picard group, usage \Pic(X)
% \Picardstack_{X/B} denotes the Picard stack of X over B
% \Picardfunctor_{X/B} denotes the Picard functor of X over B
\def\Pic{\mathop{\mathrm{Pic}}\nolimits}
\def\Picardstack{\mathcal{P}\!\mathit{ic}}
\def\Picardfunctor{\mathrm{Pic}}
\def\Deformationcategory{\mathcal{D}\!\mathit{ef}}
\usepackage{cleveref}
\crefname{originaldefinition}{original definition}{original definitions}
\Crefname{originaldefinition}{Original definition}{Original definitions}
\makeatletter
\newlabel{definition-siblings}{{0FZT}{source}{}{sourcepointer-1}{}}
\newlabel{definition-siblings@cref}{{[originaldefinition][0][]0FZT}{[1][0][]source}{}{}{}}
\makeatother
\crefname{originalreference}{original reference}{original references}
\Crefname{originalreference}{Original reference}{Original references}
\makeatletter
\newlabel{derived-lemma-representable-homological}{{0149}{source}{}{sourcepointer-2}{}}
\newlabel{derived-lemma-representable-homological@cref}{{[originalreference][0][]0149}{[1][0][]source}{}{}{}}
\makeatother
\crefname{originalreference}{original reference}{original references}
\Crefname{originalreference}{Original reference}{Original references}
\makeatletter
\newlabel{derived-lemma-split}{{05QT}{source}{}{sourcepointer-3}{}}
\newlabel{derived-lemma-split@cref}{{[originalreference][0][]05QT}{[1][0][]source}{}{}{}}
\makeatother
\crefname{originalreference}{original reference}{original references}
\Crefname{originalreference}{Original reference}{Original references}
\makeatletter
\newlabel{derived-section-ext}{{06XP}{source}{}{sourcepointer-4}{}}
\newlabel{derived-section-ext@cref}{{[originalreference][0][]06XP}{[1][0][]source}{}{}{}}
\makeatother
\crefname{originalreference}{original reference}{original references}
\Crefname{originalreference}{Original reference}{Original references}
\makeatletter
\newlabel{homology-lemma-five-lemma}{{05QB}{source}{}{sourcepointer-5}{}}
\newlabel{homology-lemma-five-lemma@cref}{{[originalreference][0][]05QB}{[1][0][]source}{}{}{}}
\makeatother
\title{Essential surjectivity and full faithfulness transfer between siblings}
\author{The Stacks Project Authors}
\date{stacks-project:equiv.tex:local-snapshot-20261012; GNU FDL 1.2 or later}
\begin{document}
\maketitle
\noindent\textit{Editorial note.} One named auxiliary 'details omitted' (five-lemma Hom check) transcribed as-is; the main Ext induction is complete in-span.\par\medskip
\noindent\textit{Editorial note.} Steps the author leaves to the reader are flagged here for the reviewer: source line(s) 2291 carry an "omitted"/"follows easily"/"immediate from the definitions" phrase. Those sentences are retained verbatim; the displayed argument around them stays inside the counted span.\par\medskip
\noindent\textit{Allowed standard prerequisites (editorial).}\begin{itemize}
\item Sibling definition and property (2). definition-siblings (same chapter, not carded). Gives the functorial isomorphisms a\_X and the factorization used throughout.
\item Five-lemma and representable-homological. homology-lemma-five-lemma, derived-lemma-representable-homological (cross-chapter, named). Feeds the auxiliary Hom-isomorphism extension step.
\item Split and Yoneda extension machinery. derived-lemma-split, derived-section-ext (cross-chapter, named). Split detection and the q>1 Yoneda composition decomposition.
\end{itemize}\medskip
\ifdefined\hypertarget\hypertarget{sourcepointer-1}{}\fi
\noindent\textit{Source reference.} Original definition 0FZT: sources/stacks-equiv/pinned-original-source.tex, source lines 2189--2189. Navigation pointer only: the counted proof cites this statement by \textbackslash{}ref as a standing input; the deduction the counted proof has to make is carried out inside the counted span, and no part of that deduction is delegated to this pointer.\par\medskip
\ifdefined\hypertarget\hypertarget{sourcepointer-2}{}\fi
\noindent\textit{Source reference.} Original reference 0149: focused-seeds/geometry/S9-eqv-sibling-faithful/ref-tags.txt, source lines 5--5. Navigation pointer only: the counted proof cites this statement by \textbackslash{}ref as a standing input; the deduction the counted proof has to make is carried out inside the counted span, and no part of that deduction is delegated to this pointer.\par\medskip
\ifdefined\hypertarget\hypertarget{sourcepointer-3}{}\fi
\noindent\textit{Source reference.} Original reference 05QT: focused-seeds/geometry/S9-eqv-sibling-faithful/ref-tags.txt, source lines 6--6. Navigation pointer only: the counted proof cites this statement by \textbackslash{}ref as a standing input; the deduction the counted proof has to make is carried out inside the counted span, and no part of that deduction is delegated to this pointer.\par\medskip
\ifdefined\hypertarget\hypertarget{sourcepointer-4}{}\fi
\noindent\textit{Source reference.} Original reference 06XP: focused-seeds/geometry/S9-eqv-sibling-faithful/ref-tags.txt, source lines 7--7. Navigation pointer only: the counted proof cites this statement by \textbackslash{}ref as a standing input; the deduction the counted proof has to make is carried out inside the counted span, and no part of that deduction is delegated to this pointer.\par\medskip
\ifdefined\hypertarget\hypertarget{sourcepointer-5}{}\fi
\noindent\textit{Source reference.} Original reference 05QB: focused-seeds/geometry/S9-eqv-sibling-faithful/ref-tags.txt, source lines 8--8. Navigation pointer only: the counted proof cites this statement by \textbackslash{}ref as a standing input; the deduction the counted proof has to make is carried out inside the counted span, and no part of that deduction is delegated to this pointer.\par\medskip
\begin{lemma}
\label{lemma-sibling-faithful}
Let $F$ and $F'$ be siblings as in Definition \ref{definition-siblings}.
Then
\begin{enumerate}
\item if $F$ is essentially surjective, then $F'$ is essentially
surjective,
\item if $F$ is fully faithful, then $F'$ is fully faithful.
\end{enumerate}
\end{lemma}

\begin{proof}
Part (1) is immediate from property (2) for siblings.

\medskip\noindent
Assume $F$ is fully faithful. Denote $\mathcal{D}' \subset \mathcal{D}$
the essential image of $F$ so that $F : D^b(\mathcal{A}) \to \mathcal{D}'$
is an equivalence. Since the functor $F'$ factors through $\mathcal{D}'$
by property (2) for siblings, we can consider the functor
$H = F^{-1} \circ F' : D^b(\mathcal{A}) \to D^b(\mathcal{A})$.
Observe that $H$ is a sibling of the identity functor.
Since it suffices to prove that $H$ is fully faithful,
we reduce to the problem discussed in the next paragraph.

\medskip\noindent
Set $\mathcal{D} = D^b(\mathcal{A})$. We have to show a sibling
$F : \mathcal{D} \to \mathcal{D}$ of the identity functor is fully faithful.
Denote $a_X : X \to F(X)$ the functorial isomorphism for
$X \in \Ob(\mathcal{A})$ given to us by Definition \ref{definition-siblings}.
For any $K$ in $\mathcal{D}$ and distinguished triangle
$K_1 \to K_2 \to K_3$ of $\mathcal{D}$
if the maps
$$
F : \Hom(K, K_i[n]) \to \Hom(F(K), F(K_i[n]))
$$
are isomorphisms for all $n \in \mathbf{Z}$ and $i = 1, 3$, then the
same is true for $i = 2$ and all $n \in \mathbf{Z}$. This uses the
$5$-lemma Homology, Lemma \ref{homology-lemma-five-lemma} and
Derived Categories, Lemma \ref{derived-lemma-representable-homological};
details omitted. Similarly, if the maps
$$
F : \Hom(K_i[n], K) \to \Hom(F(K_i[n]), F(K))
$$
are isomorphisms for all $n \in \mathbf{Z}$ and $i = 1, 3$, then the
same is true for $i = 2$ and all $n \in \mathbf{Z}$. Using the canonical
truncations and induction on the number of nonzero cohomology objects,
we see that it is enough to show
$$
F : \Ext^q(X, Y) \to \Ext^q(F(X), F(Y))
$$
is bijective for all $X, Y \in \Ob(\mathcal{A})$ and all $q \in \mathbf{Z}$.
Since $F$ is a sibling of $\text{id}$ we have $F(X) \cong X$ and
$F(Y) \cong Y$ hence the right hand side is zero for $q < 0$.
The case $q = 0$ is OK by our assumption that $F$ is a sibling of
the identity functor. It remains to prove the cases $q > 0$.

\medskip\noindent
The case $q = 1$: Injectivity. An element $\xi$ of $\Ext^1(X, Y)$
gives rise to a distinguished triangle
$$
Y \to E \to X \xrightarrow{\xi} Y[1]
$$
Observe that $E \in \Ob(\mathcal{A})$. Since $F$ is a sibling of the
identity functor we obtain a commutative diagram
$$
\xymatrix{
E \ar[d] \ar[r] & X \ar[d] \\
F(E) \ar[r] & F(X)
}
$$
whose vertical arrows are the isomorphisms $a_E$ and $a_X$.
By TR3 the distinguished triangle associated to $\xi$ we started
with is isomorphic to the distinguished triangle
$$
F(Y) \to F(E) \to F(X) \xrightarrow{F(\xi)} F(Y[1]) = F(Y)[1]
$$
Thus $\xi = 0$ if and only if $F(\xi)$ is zero, i.e., we see that
$F : \Ext^1(X, Y) \to \Ext^1(F(X), F(Y))$ is injective.

\medskip\noindent
The case $q = 1$: Surjectivity. Let $\theta$ be an element of
$\Ext^1(F(X), F(Y))$. This defines an extension of $F(X)$ by $F(Y)$
in $\mathcal{A}$ which we may write as $F(E)$
as $F$ is a sibling of the identity functor. We thus get a distinguished
triangle
$$
F(Y) \xrightarrow{F(\alpha)} F(E)
\xrightarrow{F(\beta)} F(X)
\xrightarrow{\theta} F(Y[1]) = F(Y)[1]
$$
for some morphisms $\alpha : Y \to E$ and $\beta : E \to X$.
Since $F$ is a sibling of the identity functor, the sequence
$0 \to Y \to E \to X \to 0$
is a short exact sequence in $\mathcal{A}$! Hence we obtain a
distinguished triangle
$$
Y \xrightarrow{\alpha} E \xrightarrow{\beta} X \xrightarrow{\delta} Y[1]
$$
for some morphism $\delta : X \to Y[1]$. Applying the exact functor
$F$ we obtain the distinguished triangle
$$
F(Y) \xrightarrow{F(\alpha)} F(E) \xrightarrow{F(\beta)} F(X)
\xrightarrow{F(\delta)} F(Y)[1]
$$
Arguing as above, we see that these triangles are isomorphic.
Hence there exists a commutative diagram
$$
\xymatrix{
F(X) \ar[d]^\gamma \ar[r]_{F(\delta)} & F(Y[1]) \ar[d]_\epsilon \\
F(X) \ar[r]^\theta & F(Y[1])
}
$$
for some isomorphisms $\gamma$, $\epsilon$ (we can say more but we won't
need more information). We may write $\gamma = F(\gamma')$ and
$\epsilon = F(\epsilon')$. Then we have
$\theta = F(\epsilon' \circ \delta \circ (\gamma')^{-1})$
and we see the surjectivity holds.

\medskip\noindent
The case $q > 1$: surjectivity. Using Yoneda extensions, see
Derived Categories, Section \ref{derived-section-ext}, we find that for any
element $\xi$ in $\Ext^q(F(X), F(Y))$ we can find
$F(X) = B_0, B_1, \ldots, B_{q - 1}, B_q = F(Y) \in \Ob(\mathcal{A})$ and
elements
$$
\xi_i \in \Ext^1(B_{i - 1}, B_i)
$$
such that $\xi$ is the composition $\xi_q \circ \ldots \circ \xi_1$.
Write $B_i = F(A_i)$ (of course we have $A_i = B_i$ but we don't
need to use this) so that
$$
\xi_i = F(\eta_i) \in \Ext^1(F(A_{i - 1}), F(A_i))
\quad\text{with}\quad
\eta_i \in \Ext^1(A_{i - 1}, A_i)
$$
by surjectivity for $q = 1$. Then $\eta = \eta_q \circ \ldots \circ \eta_1$
is an element of $\Ext^q(X, Y)$ with $F(\eta) = \xi$.

\medskip\noindent
The case $q > 1$: injectivity. An element $\xi$ of $\Ext^q(X, Y)$
gives rise to a distinguished triangle
$$
Y[q - 1] \to E \to X \xrightarrow{\xi} Y[q]
$$
Applying $F$ we obtain a distinguished triangle
$$
F(Y)[q - 1] \to F(E) \to F(X) \xrightarrow{F(\xi)} F(Y)[q]
$$
If $F(\xi) = 0$, then $F(E) \cong F(Y)[q - 1] \oplus F(X)$
in $\mathcal{D}$, see
Derived Categories, Lemma \ref{derived-lemma-split}.
Since $F$ is a sibling of the identity functor we have
$E \cong F(E)$ and hence
$$
E \cong F(E) \cong F(Y)[q - 1] \oplus F(X) \cong Y[q - 1] \oplus X
$$
In other words, $E$ is isomorphic to the
direct sum of its cohomology objects. This implies that the
initial distinguished triangle is split, i.e., $\xi = 0$.
\end{proof}

\end{document}
